%----------------------------------------------------------------------------------------
% Tailored CV — Akademische:r Beschäftigte:r, Institut für Pharmazeutische Biotechnologie
% Universität Ulm. Ref 26067. Deadline 18.06.2026. TV-L E13, start 01.08.2026.
% Focus: bioinformatics of high-throughput omics + AI for large datasets + teaching.
%----------------------------------------------------------------------------------------

\documentclass[10pt,a4paper,sans]{moderncv}

\usepackage[utf8]{inputenc}
\usepackage[T1]{fontenc}
\usepackage{lmodern}

\moderncvstyle{classic}
\moderncvcolor{blue}

\usepackage[scale=0.78]{geometry}
\usepackage{ragged2e}

%----------------------------------------------------------------------------------------

\firstname{Kundai Farai}
\familyname{Sachikonye}

\address{Auf der Lichtung 19}{82178 Puchheim, Germany}
\mobile{(+49) 1626386344}
\email{kundai.sachikonye@bitspark.com}
\homepage{github.com/fullscreen-triangle}
\photo[70pt][0.4pt]{pictures/Bew2.jpg}

\quote{Bioinformatician for high-throughput omics --- transcriptomics, proteomics,
and metabolomics --- with production AI tooling for large biological datasets,
and a commitment to interdisciplinary teaching.}

%----------------------------------------------------------------------------------------

\begin{document}

\makecvtitle

%----------------------------------------------------------------------------------------
%	EDUCATION
%----------------------------------------------------------------------------------------

\section{Education}

\cventry{2014--2017}{MSc Pharmaceutical Biotechnology}{HAW Hamburg}{}{}{
  Molecular biotechnology, drug development, analytical methods.
}
\cventry{2018--2019}{MSc Computational Biology}{Universität Tübingen}{}{}{
  Bioinformatics of high-throughput data: sequence analysis, variant calling,
  RNA-seq, statistical and machine learning methods for biological datasets.
}
\cventry{2019--2021}{Doctoral research, Computational Lipidomics}{Technische Universität München}{}{}{
  HPC mass-spectrometry pipelines and multi-omics integration; transitioned to
  independent research.
}
\cventry{2011--2014}{BSc Biochemistry and Cell Biology}{Jacobs University Bremen}{}{}{
  Thesis: DNA interstrand crosslinking --- mammalian cell culture and
  fluorescence microscopy.
}

%----------------------------------------------------------------------------------------
%	EXPERIENCE
%----------------------------------------------------------------------------------------

\section{Experience}

\cventry{2021--present}{Independent Bioinformatics \& Computational Biology Researcher}{Bitspark GmbH}{}{}{
  Full-time development of bioinformatic pipelines and AI tooling for
  high-throughput omics data --- genomics, transcriptomics, proteomics, and
  metabolomics. All systems documented, validated against reference data, and
  deployed as open, reproducible tools.
}

\cventry{2019--2021}{Computational Lipidomics --- Bioinformatics, HPC \& Teaching}{TUM / Bayerisches Forschungsinstitut}{Munich}{}{
  High-throughput mass-spectrometry analysis at HPC scale; multi-omics
  statistical integration; processing of $>$100\,GB datasets. Taught at
  Master's level and supervised student work.
}

\cventry{2017--2018}{Glycomics and Proteomics}{Universitätsklinikum Hamburg-Eppendorf}{Hamburg}{}{
  Automated LC-MS/MS and MALDI-TOF data pipelines; reproducible high-throughput
  workflows for clinical instrument data.
}

\cventry{2013}{Cancer Biomarker Discovery}{University Medical Center Groningen}{Groningen}{}{
  UPLC-MS/MS shotgun proteomics on tumour samples; analysis of high-dimensional
  feature sets.
}

%----------------------------------------------------------------------------------------
%	SELECTED BIOINFORMATICS SYSTEMS
%----------------------------------------------------------------------------------------

\section{Selected Bioinformatics Systems}

\cvitem{lavoisier}{
  \textbf{High-Throughput Mass-Spectrometry Pipeline.}
  Proteomics and metabolomics analysis; numerical pipeline + CNN spectral
  classification (PyTorch); distributed execution; 98.9\% NIST accuracy on
  $>$100\,GB datasets.
  \href{https://github.com/fullscreen-triangle/lavoisier}{github.com/fullscreen-triangle/lavoisier}
}

\cvitem{gospel}{
  \textbf{Multi-Omics Integration.}
  Joint modelling over host genotype, transcriptomics, and metagenomics;
  validated against 1000~Genomes, gnomAD, and UK~Biobank.
  \href{https://github.com/fullscreen-triangle/gospel}{github.com/fullscreen-triangle/gospel}
}

\cvitem{disease-state modelling}{
  \textbf{Neuroscience and Disease Characterisation.}
  Wearable and immunopeptidomic disease-state models, including
  neurodegeneration (Parkinson's, ataxia) and tumour disease state ---
  high-throughput signal and sequence analysis applied to pathological detection.
}

%----------------------------------------------------------------------------------------
%	TECHNICAL SKILLS
%----------------------------------------------------------------------------------------

\section{Technical Skills}

\cvitem{Omics bioinformatics}{
  Transcriptomics (RNA-seq), proteomics and metabolomics (mass spectrometry),
  genomics (variant calling); multi-omics integration; pathological-feature
  detection in high-throughput data
}

\cvitem{AI / ML for large datasets}{
  PyTorch (CNNs, sequence models), scikit-learn, XGBoost; Bayesian and mixture
  modelling; foundation-model fine-tuning (LoRA); reproducible AI tooling for
  biological datasets
}

\cvitem{Data engineering}{
  Ray, Dask, SLURM-compatible HPC pipelines; Docker, Kubernetes; memory-mapped
  I/O, HDF5, Parquet for $>$100\,GB datasets
}

\cvitem{Languages}{
  Python (primary), Rust, R, Bash, SQL; Nextflow / Snakemake-style workflows
}

\cvitem{Reference data}{
  1000~Genomes, gnomAD, UK~Biobank, IEDB, NIST
}

%----------------------------------------------------------------------------------------
%	TEACHING & MENTORING
%----------------------------------------------------------------------------------------

\section{Teaching \& Communication}

\cvitem{}{
  Taught at Master's level at TUM (2019--2021) and supervised student work.
  Extensive open technical documentation and tutorial material across published
  bioinformatics tools; experienced in explaining high-throughput analysis to
  interdisciplinary, non-specialist audiences. Keen to build and deliver
  bioinformatics curriculum for the new Pharmaceutical \& Molecular
  Biotechnology programme.
}

%----------------------------------------------------------------------------------------
%	LANGUAGES
%----------------------------------------------------------------------------------------

\section{Languages}

\cvitemwithcomment{English}{Native / Fluent (C2)}{}
\cvitemwithcomment{German}{Conversational}{resident in Germany since 2011}

%----------------------------------------------------------------------------------------

\renewcommand{\listitemsymbol}{-~}

\end{document}
